\section{Project Summary}

This project designed and implemented an observability-first control-plane architecture for a distributed SDN bandwidth-auction system. The resulting pipeline combines application instrumentation, centralized telemetry collection, metric-based alerting, and operator notification to support incident detection and diagnosis under realistic failure conditions.

The evaluation results indicate that cross-signal correlation across metrics, logs, traces, and alerts is effective for identifying both dependency degradation and request-path congestion. In particular, the scenario-based validation demonstrates that domain-level telemetry can expose control-plane stress conditions that are not visible through infrastructure-level health checks alone.

\subsection{Limitations}

Despite the positive evaluation outcomes, two important limitations remain.

First, system health metrics are not comprehensively monitored at the cluster level because an OpenTelemetry Operator is not currently deployed in the environment. The present workaround is the use of a custom bid-submission counter to track bid-flow behavior and support congestion detection. While this workaround is useful for the evaluated scenarios, it does not provide complete system-health coverage across all services and runtime dimensions.

Second, the current architecture places the OpenTelemetry Collector in a centralized role, which introduces a single point of failure for telemetry ingestion. If the collector becomes unavailable, metric export continuity is disrupted and observability quality degrades significantly.

\subsection{Future Directions}

A practical next step is to strengthen telemetry resilience by moving from a purely centralized collector model toward per-service sidecar collectors. Deploying an OpenTelemetry Collector sidecar alongside each service can reduce reliance on a single shared collector instance and improve fault isolation in the telemetry path.

This direction introduces a clear trade-off: improved telemetry availability at the cost of higher resource consumption and additional operational complexity. Future work should therefore evaluate this design with explicit resource-budget analysis and operational benchmarking to determine whether the resilience gains justify the overhead for production-scale deployment.

Another possible future extension is to use observed business metrics to adapt auction parameters over time. For example, if bid prices stay consistently above the reservation price for a long period, the system could raise the reservation price in a later version. That would turn the observed metrics into an input for a higher-level control policy, but it should remain a future enhancement rather than part of the current implementation.
